\chapter{Zusammenfassung}

Das Ergebnis unserer Arbeit besteht in dem funktionierendem Spiel \emph{CZGame}, welches wir als Zehnergruppe gemeinsam programmierten. 
Von der anfänglichen Spielidee wurden folgende Bereiche erfolgreich umgesetzt:

\begin{itemize}
  \item Im Spiel wurde der \textit{Pixel Art}-Stil umsesetzt.
  \item Es gibt eine Spielfigur, deren Aussehen (mittels Farbänderung) personalisiert werden kann.
  \item Der Spieler kann sich per Point and Click in einem dem realen Schulgebäude nachempfundenen Szenengebilde frei bewegen.
  \item In den Räumen der Fachschaften kann man Minigames mit verschiedenen Schwierigkeitsgraden absolvieren, wobei man bei Erfolg Items erhält.
  \item Die Spielfigur besitzt ein Inventar, in dem die Items aufbewahrt werden.
  \item Man kann den Kampf gegen einen Lehrer antreten und mithilfe der Items gewinnen.
\end{itemize}

Größtenteils wurden die zu Beginn vorgenommenen Ziele erreicht, aus zeitlichen Gründen konnten jedoch nicht alle Ideen realisiert werden. 
Folgende Bereiche der anfänglichen Spielidee konnten nicht umgesetzt werden:

\begin{itemize}
  \item In jeder Fachschaft gibt es nicht drei, sondern nur einen Lehrer.
  \item Zu Beginn des Spiels gibt es keinen Aufnahmetest, der den Spieler einer Fachschaft zuordnet, stattdessen wählt man sich selbst eine Fachschaft.
  \item Abhängig davon, welche Fachschaft man gewählt hat, bekommt man am Anfang ein (zur Fachschaft passendes) Item, später hat die gewählte Fachschaft jedoch keinen Einfluss auf den Spielverlauf.
  \item Es gibt keinen Endgegner, stattdessen hat man das Spiel gewonnen, wenn man alle Lehrer besiegt hat und sie somit in den Ruhestand geschickt hat.
\end{itemize}

Die Arbeit in der Zehnergruppe konnten wir recht gut bewältigen, jedoch brachte die Größe der Gruppe auch Herausforderungen mit sich. 
So waren in der Gruppe verschiedenste Vorkenntnisse vorhanden, was bedeutete, dass die erfahreneren Personen des Öfteren den anderen erklären mussten, 
wie man verschiedene Dinge in Java umsetzt, wodurch sie in ihrer eigenen Arbeit aufgehalten wurden. \\
Auf der anderen Seite wurden die weniger erfahrenen Personen \glqq zurückgelassen\grqq, da sie sich in kurzer Zeit in Thematiken, welche teilweise weit über dem eigenen Wissensstand liegen, einarbeiten mussten, was sehr schwer realisierbar ist.
Schlussendlich haben wir trotzdem einen guten Mittelweg gefunden, sodass jeder unabhängig vom Vorwissen seinen Teil zum Projekt beitragen konnte.
